% Canonical Paper IV appendix.
This appendix fixes matrix, action, and chart conventions used in the
projective chapters.

\subsection{One-hole arithmetic contexts}

Matrices act on the standard affine coordinate $z$ by
\[
  \begin{pmatrix}\alpha&\beta\\\gamma&\delta\end{pmatrix}\cdot z
  =\frac{\alpha z+\beta}{\gamma z+\delta}.
\]
For $c\ne0$ when required, the elementary contexts have representatives
\[
\begin{array}{c@{\qquad}c@{\qquad}c}
\toprule
context & map & matrix\\
\midrule
z+c & z\mapsto z+c & \begin{pmatrix}1&c\\0&1\end{pmatrix}\\[2mm]
z-c & z\mapsto z-c & \begin{pmatrix}1&-c\\0&1\end{pmatrix}\\[2mm]
c-z & z\mapsto c-z & \begin{pmatrix}-1&c\\0&1\end{pmatrix}\\[2mm]
cz & z\mapsto cz & \begin{pmatrix}c&0\\0&1\end{pmatrix}\\[2mm]
z/c & z\mapsto z/c & \begin{pmatrix}1&0\\0&c\end{pmatrix}\\[2mm]
c/z & z\mapsto c/z & \begin{pmatrix}0&c\\1&0\end{pmatrix}\\
\bottomrule
\end{array}
\]
Ordinary arithmetic retains the domain restriction for $c/z$ at $z=0$;
projective evaluation sends $0$ to $\infty$.  A zero multiplier or a constant
map belongs to a projective matrix monoid rather than $\PGL_2(K)$.

\subsection{Generation of the projective group}

Let
\[
  T_b(z)=z+b,
  \qquad D_a(z)=az,
  \qquad J(z)=-1/z.
\]

\begin{theorem}[Bilateral generators]
\label{thm:appendix-pgl-generation}
Translations $T_b$, nonzero scalings $D_a$, and $J$ generate
$\PGL_2(K)$.
\end{theorem}

\begin{proof}
Let
\[
  f(z)=\frac{Az+B}{Cz+D},
  \qquad AD-BC\ne0.
\]
If $C=0$, then $f$ is a nonzero scaling followed by a translation.  If
$C\ne0$, direct substitution gives
\begin{equation}
 f
 =T_{A/C}\circ
 D_{(AD-BC)/C^2}\circ
 J\circ T_{D/C}.
 \label{eq:mobius-factorization-appendix}
\end{equation}
Indeed the right side is
\[
  \frac AC-\frac{AD-BC}{C^2(z+D/C)}
  =\frac{Az+B}{Cz+D}.
\]
\end{proof}

The contexts that fix $\infty$ form the affine Borel
\[
  B_\infty=\Stab_G(\infty)
  \cong\Aff(1,K).
\]
The second-slot division $c/z$ is the elementary arithmetic context that leaves
this Borel and enters the other Bruhat cell.

\subsection{Point--predicate covariance}

For $g\in\GL(V)$, point and copoint representatives transform by
\[
  v\longmapsto gv,
  \qquad
  \varphi\longmapsto\varphi g^{-1}.
\]
Hence the pairing is invariant:
\[
  (\varphi g^{-1})(gv)=\varphi(v),
\]
and the rank-one projector is conjugated:
\begin{equation}
  \Pi_{gv,\varphi g^{-1}}
  =g\Pi_{v,\varphi}g^{-1}.
  \label{eq:projector-covariance}
\end{equation}
This explains why a bivaluation is an image--kernel result rather than two
independent scalar coordinates.

\subsection{Local affine bivaluation law}

Use the local representatives of Equation~\eqref{eq:local-projector}.  Suppose
the point coordinate obeys
\[
  a\longmapsto Aa+B,
  \qquad A\ne0.
\]
In the standard coordinate $z=-a$, a representative is
\[
  G_{A,B}=\begin{pmatrix}A&-B\\0&1\end{pmatrix}.
\]
Let $p=-1/b$ be the projective point whose line is the kernel of
$\varphi_b$.  Since $p$ transforms as a point,
\[
  p\longmapsto Ap+B.
\]
Solving again for $b$ gives
\begin{equation}
  b\longmapsto\frac{b}{A-Bb}.
  \label{eq:local-dual-affine-law}
\end{equation}
A zero denominator in Equation~\eqref{eq:local-dual-affine-law} can be a chart
escape even when the global ordered pair remains transverse.  Incidence
failure is instead $1+ab=0$.

For infinitesimal positive-real affine parameters
$A=e^{\lambda\,\mathrm{d}v}$ and $B=\mu\,\mathrm{d}u$, the first-order laws are
\[
  \mathrm{d}a=\mu\,\mathrm{d}u+\lambda a\,\mathrm{d}v,
  \qquad
  \mathrm{d}b=\mu b^2\,\mathrm{d}u-\lambda b\,\mathrm{d}v,
  \qquad
  \mathrm{d}p=\mu\,\mathrm{d}u+\lambda p\,\mathrm{d}v.
\]
These are a point/dual-point transport pair.  They are not the same binary axis
as operand-slot chirality.

\subsection{Frame coordinate above a bivaluation}

Choose the reference ordered pair $(0,\infty)$.  A third point is then any
$t\in K^\times$.  The stabilizer
\[
  H=\{z\mapsto az:a\in K^\times\}
\]
acts by $t\mapsto at$, freely and transitively.  Thus the frame fiber is a copy
of $K^\times$ only after this reference pair and affine coordinate have been
chosen.  Under a different local section, the displayed $H$-coordinate changes
by the usual gauge multiplication.

\subsection{Direct normalizer calculation}

Assume $|K|>2$ and let a projective matrix $g$ normalize the diagonal torus.
The common fixed set of the torus is $\{0,\infty\}$, so $g$ preserves this set.
If it fixes both points, it has a diagonal representative
\[
  \begin{pmatrix}a&0\\0&d\end{pmatrix};
\]
if it interchanges them, it has an antidiagonal representative
\[
  \begin{pmatrix}0&b\\c&0\end{pmatrix}.
\]
Modulo scalars these are respectively $H$ and $JH$.  This is the matrix form of
Proposition~\ref{prop:split-torus-normalizer}.

\subsection{Finite-field count from ordered frames}

An ordered projective frame over $\F_q$ is a triple of pairwise distinct points.
There are
\[
  (q+1)q(q-1)=q(q^2-1)
\]
such triples, exactly $|\PGL_2(\F_q)|$.  Forgetting the third point leaves
$(q+1)q$ bivaluations and a fiber of $q-1$ possible frame points.  This gives a
second proof of Proposition~\ref{prop:finite-field-counts} without counting
invertible matrices.
